\subsection{The complete marked local section map at every good place}

\begin{statement}{E.71}{Proved}{All ramified marked classes and every fibre of the full-to-marked local map}
Use the actual good-reduction action, oriented cusp generators and
integral tangential Frobenius lifts of E.61. Fix the marked cusp
$Q=B_0$ and put
\[
 P=\Pi,\qquad I=\overline{\langle\!\langle d_Q\rangle\!\rangle},
 \qquad W=P/\gamma_3I,\qquad B=P/I,\qquad K=I/\gamma_3I.
 \tag{E71.1}
\]
Let $v\nmid3$, $q=|k_v|$, and let $F_Q=z_QF$ be the actual
marked tangential Frobenius. Use E.67's centre
$C=\widehat{\bigwedge}^2\mathbb Z_3[[B]]$ and its exact module
\[
 Z_Q=C/(q^2\theta_Q^{[2]}-1)C,
 \qquad \theta_Q=F_Q|_B.
 \tag{E71.2}
\]
E.70 gives its finite orbit presentations and explicit choices of
coordinates identifying it with $\prod_{n\ge1}\mathbb Z_3$.

The entire ramified local $W$-section space, modulo $W$,
has the following homeomorphism:
\[
 \begin{split}
 \operatorname{Sect}_{\rm ram}(W)/W\simeq{}
 &\left((\mathbb Z_3\setminus\{0\})
       \times\mathbb Z_3/(q-1)\times Z_Q\right)\\
 &\sqcup\coprod_{i\ne Q}
       \left((\mathbb Z_3\setminus\{0\})
       \times\mathbb Z_3/(q-1)\right).
 \end{split}
 \tag{E71.3}
\]
The first summand is the actual marked-cusp family of E.67:
$t\mapsto d_Q^a t$, $F\mapsto d_Q^b c_wF_Q$.
In each of the eight other summands the representative is
\[
            t\longmapsto d_i^a t,\qquad
            F\longmapsto d_i^bF_i,
            \qquad a\ne0.
 \tag{E71.4}
\]
Every class in (E71.4) has exactly one full $P$-section class
above it. In the first summand of (E71.3), the full-source image
is exactly $w=0$, with exactly one full class above each such
point. All other points have empty full fibre.

The unramified $W$-classes and every full fibre over them are
exactly E.68.3 and E.68.6 for $R=\gamma_3I$. Thus these results
compute every local $W$-class and every fibre of
\[
          \operatorname{Sect}(P)/P
                 \longrightarrow\operatorname{Sect}(W)/W
 \tag{E71.5}
\]
at every $K$-place away from three. Formula (E71.3) describes
the ramified subspace; the unramified subspace retains its
subspace topology in the complete section space.

There is a stronger framed assertion underlying (E71.4).
Fix an actual $B$-cusp section at $i\ne Q$, including its
representative and its nonzero exponent $a$. Its actual lift
to $W$ in (E71.4) makes $K$ a group with a continuous local
action. Every $K$-valued cocycle for that action is conjugate
to the zero cocycle by exactly one element of $K$.
The comparison uses the actual path term from $Q$ to $P_i$;
it applies to every nonzero $a\in\mathbb Z_3$, including
nonunits, and to every $b\in\mathbb Z_3$.

\emph{Manuscript: \S51.16, folios733v--736v, scans172--178,
formulas(371)--(379), Part E(c),(d). The wild place, the
global compatibility problem and the other primary and
mixed finite quotients retain their stated source maps.}
\end{statement}

\begin{proof}
\emph{The actual filled curve and its ramified source.}
Let $X=C_{k_v}\setminus\{P_i:i\ne Q\}$; in this paragraph
$C_{k_v}$ denotes the smooth compact Fermat cubic.
The geometric group of $X$ is $B$, free pro-three of rank
nine. The map $P\to B$ is precisely the filling map for
the marked cusp, by E.41 and the specialization in E.61.
It carries all other peripheral generators to their
actual positive peripheral generators, denoted $c_i$.
Its base point at $Q$ becomes an ordinary rational point.

The ramified-source argument of E.61 applies to $X$ with
the following same-property comparison, which we spell out.
Its eight rational cusp vectors have the single diagonal
relation, and any two are independent. Eliminating a third
cusp in the genus-one presentation includes any chosen
pair in a free basis. The free-factor centralizer and
intersection statements used in E.61 therefore hold for
these eight inertia subgroups in $B$
\cite[Theorems A' and B']{A_HerfortRibes1985}.

For a nonzero inertia element $\nu\in B$ normalized by
$\psi=zF$, put $J=\overline{\langle\nu\rangle}$.
Suppose that $J$ lies in no conjugate of those cusp groups.
The compact-separation construction of E.61 chooses a
characteristic open $N_0\subset B$ for which the image
of $J\setminus J^3$ is disjoint from every cusp image.
The actual subgroup $H=N_0J$ is $\psi$-stable, and
\[
 H\longrightarrow J/(J\cap N_0)
                  \longrightarrow J/J^3\simeq\mathbb F_3
 \tag{E71.6}
\]
is nonzero on $J$ and kills every cusp subgroup of $H$.
The arithmetic subgroup defined by $H$ and $\psi$ gives
an actual finite cover $X'\to X$. Filling its cusps gives
the equivariant map
$H\to T_3\operatorname{Jac}(\overline X')$ through which
(E71.6) factors. Thus $\nu$ has nonzero image in this
free integral Tate module. The relation
$\psi(\nu)=\nu^q$ makes that image a $q$-eigenvector.
Every eigenvalue of this compact Jacobian Frobenius has
complex absolute value $\sqrt q$, by the same
specialization and Tate-module maps already checked in
E.61 \cite[II, Theorem1.1(b); III, Proposition9.1]{E_MilneAV2008}.
This is a contradiction. This proves the required
cuspidality for this actual filled curve.

Its free-factor centralizer is
$C_B(c_i^a)=\langle c_i\rangle$ for every $a\ne0$.
Comparing with the actual lift $F_i$ now gives exactly
$c_i^bF_i$, and conjugating by $c_i^k$ changes $b$ by
$(1-q)k$. Distinct cusps and distinct nonzero exponents
remain distinct in the integral abelianization.
This proves that the ramified $B$-section space consists
of its eight actual cusp families. Each has the displayed
full $P$ lift, so the comparison needed below is onto.

\emph{The exact affine action on the marked kernel.}
Set $D=I/I'=\mathbb Z_3[[B]]$; E.67 supplies the actual
universal maps and the compact central sequence
\[
                  0\longrightarrow C\longrightarrow K
                         \longrightarrow D\longrightarrow0.
 \tag{E71.7}
\]
For $i\ne Q$, put
$b_i=\overline{z_i z_Q^{-1}}\in B$.
Since $F_i=(z_i z_Q^{-1})F_Q$, the actual action of
$F_i$ on $D$ is $q$ times the affine permutation
$g\mapsto b_i\theta_Q(g)$.
For the reference section (E71.4), its inertia action
on $D$ is the permutation $T_0:g\mapsto c_i^a g$, and
its Frobenius action is $qT$, where
\[
               T(g)=c_i^b b_i\theta_Q(g).
 \tag{E71.8}
\]
Its actions on $C$ are the corresponding diagonal
permutation $T_0^{[2]}$ and $q^2T^{[2]}$.
These formulas follow directly on the generators
$g d_Qg^{-1}$ of $I$ and on their commutators.
Possible central terms in the action on $K$ are
retained by (E71.7); they do not change these two
induced actions.

Put $H_a=\langle c_i^a\rangle$. The relation
$T T_0T^{-1}=T_0^q$ shows that $T$ acts on the compact
left coset space $H_a\backslash B$ and on
$H_a\backslash(B\times B)$, with diagonal left
translation in the latter. The closure of this
permutation action is pro-three: it comes from the
actual pro-three semidirect product of the geometric
group with the local arithmetic quotient of E.61.
All its finite permutation orbits therefore have
three-power length.

\emph{No periodic coset: the actual compact Kummer map.}
Use $Q$ as origin on the compact elliptic curve $E$.
Filling all the remaining cusps gives the actual
equivariant surjection
\[
                 p:B\longrightarrow T_3E .
 \tag{E71.9}
\]
Every $p(c_i)$ is zero. The nine original coordinate
cusps are the nine distinct three-torsion points on
$E$, also after good specialization. Indeed each
coordinate cusp is a flex, whose tangent cuts out
three times that point. Comparing with the tangent
at $Q$ gives $3(P_i-Q)=0$, and distinct points on
the genus-one curve have distinct Abel--Jacobi images.
The disjoint integral sections of E.61 preserve this
nonzero three-torsion at every residue characteristic
other than three.

With the path direction used for $F_i=z_iF$, the
cocycle $p(b_i)$ is the Kummer cocycle of $P_i-Q$.
This assertion follows on every actual division
cover $[3^n]:E\to E$: transporting the zero fibre
from $Q$ to $P_i$ chooses a compatible division point
$R_n$ with $[3^n]R_n=P_i-Q$, and the Frobenius
translation is $F_E(R_n)-R_n$. Changing the chosen
path adds the corresponding Tate-module coboundary.
Thus the comparison retains both its actual
representative and its precisely specified ambiguity.

Suppose $H_ag$ were periodic for $T$, of period $m>0$.
Then $T^m(g)=c_i^{ak}g$ for some $k\in\mathbb Z_3$.
Projecting by (E71.9) gives
\[
 \sum_{j=0}^{m-1}F_E^j p(b_i)
                         =(1-F_E^m)p(g).
 \tag{E71.10}
\]
The terms from $c_i^b$ vanish under $p$.
Equation (E71.10) would make the restriction of that
actual Kummer cocycle to $k_{v,m}=\mathbb F_{q^m}$
a coboundary. More precisely, (E71.10) makes
$R_n+p(g)_n$ fixed by $F_E^m$, where $p(g)_n$
is the $3^n$-torsion coordinate of the actual Tate
element $p(g)$. These are compatible
$k_{v,m}$-rational divisions of $P_i-Q$
at every level. The finite group $E(k_{v,m})$ has
no nonzero three-primary element divisible by
arbitrarily high powers of three. Since $P_i-Q$
is nonzero of order three, this is impossible.
Hence $T$ has no periodic point on $H_a\backslash B$.
This proof includes every nonunit $a$ and every $b$;
it uses the actual affine term $b_i$ in (E71.8).

\emph{The compact measure and coinvariant maps.}
For any continuous permutation with pro-three
closure on a compact profinite space and with no
periodic points, its invariant integral measures
are zero. Here is the exact trace argument needed.
Choose a cofinal system of invariant finite
partitions. Bounded orbit lengths along all
refinements would, by compactness, produce a
point fixed by some power $T^{3^e}$. Thus the
minimum orbit length tends to infinity. At a
fixed coarse orbit, refining to orbit lengths
divisible by $3^n$ times its length makes the
coarse coefficient of every invariant measure
divisible by $3^n$. Every such coefficient is
therefore zero.

The actual augmentation to cosets gives
\[
 \begin{split}
 D/(T_0-1)D&\simeq\mathbb Z_3[[H_a\backslash B]],\\
 C/(T_0^{[2]}-1)C
  &\simeq e^-\mathbb Z_3[[H_a\backslash(B\times B)]],
       \qquad e^-=(1-\mathrm{swap})/2.
 \end{split}
 \tag{E71.11}
\]
For the first map, $\delta_g$ goes to $\delta_{H_ag}$.
For the second, use the alternating-measure
embedding of E.67, then pass to diagonal orbits.
The idempotent $e^-$ commutes with diagonal
translation and with all quotient maps. At a
finite quotient these are the ordinary exact
coinvariant maps of finite permutation modules.
Passing to the compact inverse limits is exact:
the finite systems of lift equations have the
finite intersection property. The images of
$T_0-1$ and $T_0^{[2]}-1$ are already compact
and closed. This proves (E71.11), without choosing
a coset transversal.

The no-periodic-point result and the trace
argument show
\[
     \left(D/(T_0-1)D\right)^T=0 .
 \tag{E71.12}
\]
Also $D^{T_0}=0$ and $C^{T_0^{[2]}}=0$.
For these two assertions, left translation by
$c_i^a$ acts freely on $B$ and diagonally on
$B\times B$, and its finite orbit lengths tend
to infinity. The same trace argument applies;
$C$ is an invariant direct summand of the
second measure module. The assertions also
hold for $T_0^q$, which generates the same
closed translation group.

Finally, on the second coinvariant module of
(E71.11) the operator $qT^{[2]}-1$ is injective.
In each finite integral orbit module $T^{[2]}$
is a permutation $P$ of finite order $N$. A
vector in its kernel is killed by $q^N-1$.
That nonzero integer acts injectively on the
free $\mathbb Z_3$-module, and hence on its
alternating direct summand. The compact limit
therefore has zero kernel as well.

\emph{Eliminating the entire cocycle, including its central part.}
Let a $K$-valued cocycle for the reference
local action have inertia and Frobenius values
whose images in $D$ are $x,u$. The tame
relation gives the exact additive equation
\[
       \bigl(qT-N_q(T_0)\bigr)x
                            =(T_0^q-1)u,
       \qquad N_q(T_0)=1+T_0+\cdots+T_0^{q-1}.
 \tag{E71.13}
\]
Modulo $(T_0-1)D$ this reads $q(T-1)\bar x=0$.
Since $q$ is a three-adic unit, (E71.12) gives
$x=(T_0-1)v$ for one and only one $v\in D$.
Lift $v$ to an actual element of $K$ and
conjugate the section by it. This replaces
$x$ by $x-(T_0-1)v=0$.
The new equation gives $(T_0^q-1)u=0$, so
the new Frobenius image in $D$ is zero.
The entire new cocycle consequently takes
values in $C$.

Write these central values as $x_C,u_C$.
Their exact relation is
\[
 \bigl(q^2T^{[2]}-N_q(T_0^{[2]})\bigr)x_C
                       =((T_0^{[2]})^q-1)u_C.
 \tag{E71.14}
\]
On its inertia coinvariant module the equation
is $q(qT^{[2]}-1)\bar x_C=0$.
The proved injectivity gives
$x_C=(T_0^{[2]}-1)v_C$ for a unique $v_C\in C$.
Conjugating by $v_C$ kills this inertia value;
the remaining Frobenius value is zero because
$C^{(T_0^{[2]})^q}=0$.
Thus every cocycle is conjugate to zero.
An element fixing the zero cocycle has
$T_0$-fixed image in $D$, hence has image
zero there; its remaining central value is
also $T_0^{[2]}$-fixed and is zero.
This proves uniqueness of the full conjugator.
All higher central corrections in the original
action on $K$ have been handled by the actual
first conjugation and the subsequent central
equation; none was replaced by a graded action
on the whole group.

\emph{Every geometric orbit and every full fibre.}
Take a $W$-section whose inertia has nontrivial
image in $B$. Normalize its projected section
to its unique cusp parameters $(i,a,\bar b)$.
Lift its conjugator through $W\to B$.
The remaining section is a $K$-valued cocycle
over the actual reference (E71.4), so the
preceding argument conjugates it to that
reference. Conversely two references
equivalent in $W$ have equivalent $B$ images.
The $B$ classification forces the same cusp,
the same exponent, and exactly the stated
change of $b$; that change is realized in
$W$ by the actual lift $d_i^k$.

The other ramified $W$-sections have inertia
image one in $B$, so E.67 gives exactly the
first summand of (E71.3), including its
full-source image. The map $P\to W$ preserves
the full abelianization. E.61 therefore shows
that an unramified target can have only
unramified full preimages, and that every
class in (E71.4) has precisely its displayed
full cusp class above it. This proves all
set-theoretic assertions and the framed
assertion.

For topology, each fixed nonzero valuation
annulus for $a$ has compact parameter space,
and its actual representative map to the
Hausdorff geometric orbit space is continuous
and bijective. It is therefore a homeomorphism.
Such annuli are open in and cover the ramified
space, as detected by the valuation of its nonzero
abelian inertia vector. The
nonzero cusp vectors in the free integral
abelianization distinguish its nine branches;
their intersections occur only at zero,
outside this ramified space. This gives the
homeomorphism (E71.3). The unramified topology
and all its full fibres are those already
proved with the actual unique conjugators
in E.68. The retained program independently enumerates
the entire projective cubic over $\mathbb F_4$ and
checks the literal tangent polynomial at all nine flexes.
Using the actual compact Frobenius $-2$ of E.37, it
also checks the Kummer obstruction for every one of
the eight nonzero three-torsion coordinates over
all residue extension degrees from one through 81.
All 648 nonzero residues lie outside the actual
Frobenius coboundary image; the argument above
proves that assertion for every positive degree.
This completes (E71.5).
\end{proof}
